% Copyright (c) 2026 Eugene Joseph M. Ragasa.
% SPDX-License-Identifier: Apache-2.0
% Author: Eugene Joseph M. Ragasa
% Formal proof completed: 2026-10-19
\documentclass[11pt]{article}
\usepackage[T1]{fontenc}
\usepackage{amsmath,amssymb,booktabs}
\usepackage[margin=1in]{geometry}
\usepackage[hidelinks]{hyperref}
\usepackage{tabularx}
\title{Proof Dossier: Ragasa2026A Aligned Fields Have Common Power}
\author{Eugene Joseph M. Ragasa}
\date{Formal proof completed: 2026-10-19}
\begin{document}
\maketitle
\section{Identity and status}
\begin{tabularx}{\textwidth}{@{}l>{\raggedright\arraybackslash}X@{}}
\toprule
Statement ID & \texttt{BG-RAGASA2026A-THM-002} \\
Status & \texttt{proved} \\
Lean declaration & \path{ProjectKoios.BandGauge.Papers.Ragasa2026A.alignedSplittingFields_commonPower} \\
Module & \path{ProjectKoios/BandGauge/Papers/Ragasa2026A/RankTwoFields.lean} \\
\bottomrule
\end{tabularx}

\section{Exact statement}
If the centered sampled parent has source-normalized weighted Frobenius power $\kappa$,
then
\[
 P(D_{\phi})=\kappa,\qquad P(B)=\kappa.
\]
The contract retains the source premises $\alpha_i>0$, $\sum_i\alpha_i=1$, and
$E_{\mathrm{ref}}>0$.

\section{Source and proof}
This theorem \texttt{formalizes} the common-power consequence of manuscript lines
603--624 at immutable \path{ksdft2effmass} commit
\begin{center}\small\texttt{99a46cb59d0c6f27ada7d709f194996ef9619d9f}.\end{center}
The global standard real rotation and every supplied comparison alignment are unitary, so
weighted Frobenius power is preserved pointwise and then under the finite sum.

\section{Boundary and trusted base}
The parent-power equality is a hypothesis. No concrete matrix or numerical power is
validated. Lean and pinned Mathlib report \texttt{propext},
\texttt{Classical.choice}, and \texttt{Quot.sound}. No admission or scientific axiom is
used. AI-assisted and kernel-checked; no independent expert review is recorded. Status:
\texttt{proved}; working draft.
\end{document}
